\chapter{Baker's Cysts}
\index{Baker's Cysts}

\section*{Definition and Overview}
A Baker's cyst (also called a popliteal cyst) is a fluid-filled swelling at the back of the knee. Despite the name, it is not a true cyst but an outpouching of the knee joint's lining: extra joint fluid is pushed backwards through a small opening into a bursa (a fluid sac) behind the knee, creating a bulge. It is named after the surgeon William Morrant Baker, who described it in the 1800s.

Baker's cysts are almost always a sign of another problem in the knee that is producing too much fluid, such as osteoarthritis, rheumatoid arthritis, or a meniscal tear. Many cause no symptoms and are found incidentally. The main concern is that a sudden calf swelling could be a ruptured cyst, which looks similar to a blood clot (deep vein thrombosis) and must be checked.

\section*{Epidemiology}
Baker's cysts are common, especially in adults over 40, in whom knee joint problems such as osteoarthritis and rheumatoid arthritis are frequent. They occur at all ages, including children, where they often appear without any underlying knee condition and usually disappear on their own. Cysts are often found incidentally on imaging of the knee, and are present in a substantial proportion of people having an MRI or ultrasound for other reasons. They are more likely to become symptomatic when the knee joint is actively inflamed.

\section*{Aetiology and Pathophysiology}
A Baker's cyst forms when fluid builds up inside the knee joint (an effusion) and is forced through a weak spot in the joint capsule into the bursa at the back of the knee. Common underlying causes include:

\begin{itemize}
    \item \textbf{Osteoarthritis} -- the most common cause in older adults
    \item \textbf{Rheumatoid arthritis} and other inflammatory joint conditions
    \item \textbf{Meniscal tears} and cartilage damage
    \item \textbf{Ligament injuries}, such as an ACL tear
    \item \textbf{Gout or pseudogout}, which cause joint inflammation
    \item \textbf{In children:} often no cause is found, and the cyst resolves on its own
    \item Rarely, \textbf{joint infection} (septic arthritis) can cause a painful cyst and needs urgent treatment
\end{itemize}

\section*{Clinical Features}
\begin{itemize}
    \item A swelling or bulge at the back of the knee that may be felt as a soft lump
    \item A feeling of tightness or fullness behind the knee
    \item Discomfort or pain, especially when bending the knee, kneeling, or during exercise
    \item In many cases, no symptoms at all -- the cyst is noticed on imaging
    \item If the cyst is large: stiffness or reduced range of movement
    \item If it ruptures: sudden pain and swelling of the calf with bruising, closely resembling a deep vein thrombosis (DVT)
\end{itemize}

\section*{Differential Diagnosis}
\begin{itemize}
    \item \textbf{Deep vein thrombosis (DVT):} a blood clot in the calf; the most important condition to exclude when the calf swells
    \item \textbf{Popliteal aneurysm:} a bulging artery behind the knee, felt as a pulsatile swelling
    \item \textbf{Muscle injury or tear} of the calf or hamstring
    \item \textbf{Soft tissue tumour} around the knee (rare)
    \item \textbf{Enlarged lymph nodes} behind the knee
    \item \textbf{Other knee swellings,} such as a meniscal cyst or ganglion
\end{itemize}

\section*{When to Seek Medical Care}
Any new swelling, pain, or redness of the calf should be checked promptly, because it can be difficult to tell a ruptured Baker's cyst from a deep vein thrombosis without tests.

Seek medical assessment if:
\begin{itemize}
    \item You notice a new lump or swelling behind the knee
    \item The calf suddenly becomes painful, swollen, warm, or red
    \item The swelling is associated with knee pain, locking, or giving way
    \item The lump is rapidly enlarging or causing numbness, tingling, or weakness
\end{itemize}

\textbf{Contact your local emergency services if you develop sudden chest pain or difficulty breathing} together with calf swelling -- this could suggest a blood clot that has travelled to the lungs.

\section*{Diagnosis and Investigations}
\begin{itemize}
    \item \textbf{History and examination:} a bulge behind the knee that is more obvious with the knee extended; the knee itself is examined for fluid, tenderness, and stability
    \item \textbf{Ultrasound:} the key test; it confirms the cyst, shows whether it has ruptured, and distinguishes it from a DVT or popliteal aneurysm
    \item \textbf{Knee X-ray:} used to assess for osteoarthritis
    \item \textbf{MRI:} provides detailed images if an underlying meniscal tear or ligament injury needs definition
    \item \textbf{Blood tests:} inflammatory markers or uric acid may help if arthritis or gout is suspected
    \item \textbf{Aspiration:} fluid may be drawn off the cyst, rarely, to exclude infection or to relieve symptoms
\end{itemize}

\section*{Management}
\begin{itemize}
    \item \textbf{Treat the underlying knee problem:} managing osteoarthritis or rheumatoid arthritis is the most important step, because the cyst usually shrinks as joint fluid settles
    \item \textbf{Physiotherapy and exercise:} strengthening and stretching programs for the knee can reduce fluid production and symptoms
    \item \textbf{Pain relief:} simple analgesics and anti-inflammatory medicines for knee symptoms, and measures to reduce joint swelling
    \item \textbf{Rest, ice, and elevation:} helpful during a flare-up of knee swelling
    \item \textbf{Injection:} aspiration of the cyst combined with a steroid injection into the knee may be offered for large, painful cysts
    \item \textbf{Surgery:} arthroscopic treatment of an underlying meniscal tear may be considered; removal of the cyst itself is rarely needed
    \item \textbf{Children:} usually only reassurance and observation, since these cysts typically resolve on their own
\end{itemize}

\section*{Complications}
\begin{itemize}
    \item Rupture of the cyst, causing sudden calf pain, swelling, and bruising that mimics a DVT
    \item Rarely, the cyst compresses nearby structures, causing numbness, tingling, or leg swelling
    \item Pain and restricted movement if the cyst is large
    \item Recurrence after aspiration is common if the underlying knee condition continues
    \item Rarely, infection of the cyst, which can occur with a septic knee joint
    \item Missing an underlying problem, such as a meniscal tear, that needs separate treatment
\end{itemize}

\section*{Prognosis and Outcomes}
Baker's cysts are benign and not cancerous, and they do not spread to other parts of the body. The outlook depends mainly on the underlying knee condition: when the knee problem is treated and joint swelling settles, the cyst often shrinks or disappears. Some cysts persist but remain harmless and symptom-free. In children, most cysts resolve spontaneously over months to years. A ruptured cyst settles over several weeks with rest and elevation, but requires checking first to rule out a blood clot. Surgical removal is rarely needed, and recurrence after any treatment is possible.

\section*{Prevention and Self-Care}
\begin{itemize}
    \item Keep the knee joint healthy: maintain a healthy weight and stay active with low-impact exercise
    \item Strengthen the thigh (quadriceps) and hamstring muscles to support the knee
    \item Seek treatment for knee osteoarthritis, rheumatoid arthritis, and other joint conditions to reduce fluid buildup
    \item After activity, rest and elevate the leg if the knee feels swollen
    \item Ice the knee and use simple pain relief during flares
    \item If the calf suddenly swells or becomes painful, seek assessment promptly to exclude a blood clot
    \item Follow up with your doctor if a cyst is large, painful, or not settling
\end{itemize}

\section*{Sources and Further Reading}
The following sources provide reliable, up-to-date information on Baker's cysts.
\begin{itemize}
    \item NHS. \textit{Baker's cyst.} https://www.nhs.uk/conditions/
    \item Mayo Clinic. \textit{Baker's cyst.} https://www.mayoclinic.org
    \item MSD Manual. \textit{Baker Cyst (Popliteal Cyst).} https://www.msdmanuals.com
    \item MedlinePlus. \textit{Baker's cyst.} https://medlineplus.gov
\end{itemize}
